\begin{textsample}{NEG Dim 2 – global\_north\_2024\_04 – Score 1.00 – global\_north\_2024\_04/200415979.txt}  \label{ex:f2_neg_185}
Angle down iconAn icon in the shape of an angle pointing down.Israeli soldiers of the Netzah Yehuda Haredi infantry battalion hold their weapons during their swearing-in ceremony in Jerusalem on May 26, 2013.
REUTERS The US is set to sanction an IDF unit following reported human rights abuses, Axios reports.
The Netzah Yehuda battalion was implicated in the death of an 80-year-old US citizen in 2022.
The news of the possible sanctioning of the IDF unit came as the House of Representatives on Saturday passed a bill that includes more than \$14 billion in military aid to Israel as well as \$9 billion in humanitarian assistance, much of which will go to Gaza.
Advertisement The unprecedented sanctions will prohibit the battalion and its members from receiving US military assistance or training.
The imposed measures are in accordance with the 1997 Leahy vetting policy, which forbids US foreign aid and Defense Department training programs from going to foreign security, military, and police units credibly alleged to have committed human rights violations.
In 2022, the State Department requested the US @ @ @ @ @ @ @ @ @ @ in the West Bank following reports of abuse of Palestinians, the Israeli outlet Haaretz, reported.
Last month, a State Department panel recommended that Blinken disqualify multiple Israeli military and police units operating in the West Bank from receiving US aid, ProPublica reported.
Advertisement Blinken confirmed on Friday that he had made determinations based on this investigation and stated,"You can expect to see them in the days ahead"regarding the sanctions, per Axios.
Related story Senior members of the Israeli government pushed back at the potential sanctions move against the Netzah Yehuda battalion.
Radical settlers Israeli settlers hold a protest march from Tapuach Junction in the Israeli-occupied West Bank, April 10, 2023.
REUTERS/Nir Elias The Netzah Yehuda unit, part of the Kfir brigade, was originally set up in 1999 to accommodate the religious beliefs of recruits from ultra-Orthodox and national religious communities, including those from extremist settlements.
All of its members are men.
The unit has a reputation for recruiting some @ @ @ @ @ @ @ @ @ @, said Axios.
Historically, the Netzah Yehuda battalion has been primarily deployed in the West Bank, gaining notoriety for its treatment of Palestinians.
Advertisement Notably, soldiers from the unit were accused in the death of Omar Assad, an 80-year-old US citizen, who died of a heart attack in 2022 after being detained, bound, gagged, and abandoned by members of the battalion.
While Israel has deployed the Netzah Yehuda battalion out of the West Bank since the incident, the unit has continued to serve -- primarily in the country ' s north and in the Gaza Strip during the ongoing conflict with Hamas.

% matched lemmas: 
\end{textsample}
